# TA-BRPL 디펜스 현장 참고 시트

## ★ 디펜스 황금 원칙
> 한계 인정 → **"그럼에도 기여는 X다"** 반드시 세트로

---

## 예상 질문 & 즉답 카드

---

### [Q1] "att_share 1.7%p가 실용적으로 의미 있는가?"

**즉답:**
"route capture의 제거가 아닌 감소이며, 서브트리 구조에서 att_share 감소는
하위 노드 전체의 피해 완화로 이어집니다.
75쌍 전체에서 76% 승률로 일관되게 관찰된 것이 핵심입니다.
standalone 완전 방어가 아닌 경량 완화책으로서의 기여를 주장합니다."

**원고 위치:**
- Abstract: `lightweight trust penalty` → `on-node trust penalty` 수정 예정
- 결론 절: `att_share` 해석 문장
- Table 3: 절대값 비교 (RPL=0.110, BRPL=0.110, TA-BRPL=0.093)
- Table 4: H1 p < 10⁻⁴, 95% CI [-0.023, -0.011]

---

### [Q2] "lightweight가 검증되지 않은 주장 아닌가?"

**즉답:**
"맞습니다. 현재 원고에서 lightweight는 실측 비용 claim이 아니라
PKI/IDS 없이 로컬에서 처리한다는 설계 의도입니다.
TelosB에서 pow()와 double 연산 비용은 실측하지 않았으며,
원고 표현을 on-node, cryptography-free로 수정할 예정입니다."

**원고 위치:**
- Abstract 첫 문단: "lightweight trust penalty"
- 결론 절: "lightweight defense"
- 7절 한계: 에너지/지연 미측정 명시 (이미 존재)

---

### [Q3] "가중치 0.5/0.3/0.2 근거가 없다"

**즉답:**
"이론적 최적값이 아닌 경험적 operating point입니다.
고정 토폴로지에서 1차 안정화 후 랜덤 75쌍에서 재현했고,
ablation에서 세 신호 결합의 필요성은 별도로 검증했습니다.
sensitivity analysis 부재는 한계로 인정하며 향후 과제입니다."

**원고 위치:**
- 4.2절: `(wf, wc, wh) = (0.5, 0.3, 0.2)` 설명
- Table 2: 설계 버전별 성능 격차 (v1~v13.12)
- Table 6: Ablation 결과 (Full vs Tfwd only vs Tfwd+Tctrl)
- 7절: 파라미터 민감도 한계

---

### [Q4] "파라미터 튜닝셋과 평가셋이 분리됐는가?"

**즉답:**
"고정 토폴로지 튜닝과 랜덤 75쌍은 1차 분리됩니다.
그러나 ablation에 쓴 일부 토폴로지/시드가 메인셋에 포함되어
strict hold-out은 아닙니다.
따라서 '부분적으로 겹치는 corpus에서 재현했다'는 표현이 정확하며,
held-out generalization 주장은 원고에서 삭제할 것입니다."

**원고 위치:**
- 5.1절: 실험 행렬 설명
- 7절: 파라미터 민감도 항목
- Table 2: 고정 토폴로지 기준 버전별 결과

---

### [Q5] "혼잡 이점 보존을 검증했는가?"

**즉답:**
"직접 검증하지 않았습니다.
혼잡 유도 구간을 포함했으나 load balancing 보존 여부를
별도 지표로 평가하지 않았습니다.
정확한 주장은 '혼잡 유도 workload를 포함한 공격 실험에서
PDR 비열세를 보였다'입니다."

**원고 위치:**
- 5.1절: 시뮬레이션 시간 설정 (워밍업/공격/복구)
- 5.2절: 평가 지표 (att_share, hit_ratio, PDR_dur, churn)
- 7절: 에너지/지연 오버헤드 항목

---

### [Q6] "표 1 비교만 하고 실험 비교가 없다"

**즉답:**
"표 1은 문헌 기반 설계 위치 비교이며 성능 비교표가 아닙니다.
따라서 '설계 위치가 다르다'는 주장만 방어 가능합니다.
RFTrust, PCC-RPL과의 직접 비교는 재구현 공수 문제로
향후 과제로 미뤘으며, 이는 7절에 명시되어 있습니다.
우월성 함의를 주는 문장은 원고에서 수정할 것입니다."

**원고 위치:**
- 2.6절: 표 1 비교표
- Table 1 캡션
- 7절: 공격 범위의 제한 항목 (RFTrust, PCC-RPL 언급)
- 향후 연구: 에너지·지연 오버헤드 및 동급 비교 항목

---

### [Q7] "왜 accept해야 하는가 — 핵심 기여는?"

**즉답 (5줄 버전):**

"첫째, sinkhole 평가 목표를 PDR에서 route capture 억제로 재정의했습니다.
둘째, trust를 BRPL 점수 내부에 소프트 패널티로 직접 결합하는
설계 구조를 제시했습니다.
셋째, BRPL의 queue-awareness만으로는 sinkhole을 완화할 수 없다는
부정적 실증 결과를 냈습니다 — 이것이 단순 성능 수치보다 더 중요합니다.
넷째, 다중 신호 결합 + 안정화 정책이 함께 필요하다는 설계 법칙을
ablation으로 실증했습니다.
다섯째, 밀도별 층화 75쌍 + 5시드 + 사전 선언 가설의 방법론적 틀을
후속 연구가 재사용할 수 있도록 제시했습니다."

**원고 위치:**
- 1절 서론: C1~C4 기여 목록
- Table 3: RPL=BRPL att_share (부정적 실증)
- Table 4: 가설 검증 결과
- Table 6: Ablation
- 5.1절: 실험 설계 방법론

---

## 핵심 수치 암기 카드

| 지표 | BRPL | TA-BRPL | Δ | p값 |
|------|------|---------|---|-----|
| att_share | 0.110 | 0.093 | -0.0169 | < 10⁻⁴ |
| hit_ratio | 0.117 | 0.101 | -0.0166 | < 10⁻⁴ |
| PDR_dur | 0.950 | 0.961 | +0.0102 | 0.002 |
| churn | 0.037 | 0.073 | +0.0364 | CI 상한 0.057 < 0.1 |

- 승률: att_share 76%, 중밀도 84%
- PDR 비열세 통과율: 93.3% (실패 5건 모두 저밀도)
- 사전 선언 churn 상한: +0.1

---

## 원고 수정 우선순위 체크리스트

### 즉시 수정 (디펜스 전)
- [ ] `lightweight` 표현 → `on-node` / `cryptography-free` 교체 (Abstract, 결론)
- [ ] 표 1 캡션에 "qualitative positioning" 명시
- [ ] 우월성 암시 문장 삭제 또는 완화

### 디펜스 후 revision
- [ ] 7절 한계에 tuning/evaluation 분리 불완전성 추가
- [ ] 혼잡-only 비교 미수행을 한계절에 추가
- [ ] "held-out generalization" 표현 삭제
- [ ] 가중치 선택 근거 문장 명확화
- [ ] Table 1 관련 연구 비교 한계 각주 추가